% Configuration globale pour un mémoire de soutenance professionnel
% Ce fichier définit la mise en page et les styles pour l'ensemble du document

% ====== PACKAGES ESSENTIELS ======
\usepackage{fontspec}      % Support des polices TrueType/OpenType
\usepackage{polyglossia}   % Support multilingue
\setmainlanguage{french}   % Langue principale
\setotherlanguage{english} % Langue secondaire

% ====== POLICES GOOGLE FONTS ======
\defaultfontfeatures{
    Path = ./fonts/,
    Extension = .ttf,
    Ligatures=TeX,
    Scale=1.0
}

% Police principale pour tout le document
\setmainfont{Roboto}[
    UprightFont = *-Regular,
    BoldFont = *-Bold,
    ItalicFont = *-Italic,
    BoldItalicFont = *-BoldItalic
]

% Police sans-serif (pour les titres et éléments d'interface)
\setsansfont{Roboto}[
    UprightFont = *-Regular,
    BoldFont = *-Bold,
    ItalicFont = *-Italic,
    BoldItalicFont = *-BoldItalic
]

% Police pour les mathématiques
\setmathfont{Latin Modern Math}

% ====== MISE EN PAGE ======
\usepackage{geometry}
\geometry{
  a4paper,           % Format A4
  top=2.5cm,         % Marge supérieure
  bottom=2.5cm,      % Marge inférieure
  left=2.5cm,        % Marge gauche
  right=2.5cm,       % Marge droite
  headheight=14pt,   % Hauteur de l'en-tête
  footskip=1.5cm     % Espace pour le pied de page
}

% ====== EN-TÊTES ET PIEDS DE PAGE ======
\usepackage{fancyhdr}
\pagestyle{fancy}
\fancyhf{}
\fancyhead[L]{\sffamily\nouppercase{\leftmark}}
\fancyhead[R]{\sffamily\thepage}
\fancyfoot[C]{\sffamily\small Mémoire de soutenance}
\renewcommand{\headrulewidth}{0.4pt}
\renewcommand{\footrulewidth}{0.2pt}

% ====== STYLES DE TEXTE ======
\usepackage{setspace}
\onehalfspacing  % Interligne 1.5

% ====== STYLES DES FIGURES ======
\usepackage{caption}
\usepackage{subcaption}
\captionsetup{
  font=small,
  labelfont={bf,sf},
  textfont=sf,
  format=hang,
  width=0.9\textwidth,
  skip=10pt
}

% ====== LIENS HYPERTEXTE ======
\usepackage{hyperref}
\hypersetup{
  colorlinks=true,
  linkcolor=blue,
  citecolor=blue,
  urlcolor=blue,
  bookmarks=true,
  bookmarksopen=true,
  bookmarksnumbered=true,
  pdfauthor={\theauthor},
  pdftitle={Mémoire de Soutenance}
}

% ====== STYLES DES LISTES ======
\usepackage{enumitem}
\setlist{nosep}  % Réduit l'espace entre les éléments de liste

% ====== TABLE DES MATIÈRES ======
\usepackage{tocloft}
\renewcommand{\cfttoctitlefont}{\huge\bfseries\sffamily\color{bookColor}}
\renewcommand{\cftchapfont}{\bfseries\sffamily}
\renewcommand{\cftsecfont}{\sffamily}

% ====== STYLE DES TABLEAUX ======
\usepackage{booktabs}  % Pour des tableaux professionnels
\usepackage{array}
\usepackage{tabularx}
\usepackage{colortbl}

% ====== CODE SOURCE ======
\usepackage{listings}
\lstset{
  basicstyle=\footnotesize\ttfamily,
  breaklines=true,
  frame=single,
  numbers=left,
  numberstyle=\tiny,
  keywordstyle=\color{blue},
  commentstyle=\color{green!60!black},
  stringstyle=\color{red}
}
